\begin{titlepage}
  \begin{center}

  {\Huge BUS\_1553}

  \vspace{25mm}

  \includegraphics[width=0.90\textwidth,height=\textheight,keepaspectratio]{img/AFRL.png}

  \vspace{25mm}

  \today

  \vspace{15mm}

  {\Large Jay Convertino}

  \end{center}
\end{titlepage}

\tableofcontents

\newpage

\section{Usage}

\subsection{Introduction}

\par
BUS1553 is a core for interfacing the PMOD1553 device to a bus of choice. The core will process data to and from the PMOD1553.
The data can then be accessed over a BUS, currently AXI lite or Wishbone Standard, and processed as needed. All input and output
over the bus goes into FIFOs that is then tied to the demodulation and modulation cores, which then send/recv the differential data
to/from the PMOD1553 device. The following is information on how to use the device in an FPGA, software, and in simulation.

\subsection{Dependencies}

\par
The following are the dependencies of the cores.

\begin{itemize}
  \item fusesoc 2.X
  \item iverilog (simulation)
  \item cocotb (simulation)
\end{itemize}

\subsubsection{axi\_lite\_1553 Depenecies}
\begin{itemize}
\item dep
	\begin{itemize}
	\item AFRL:utility:helper:1.0.0
	\item AFRL:device:up\_1553:1.0.0
	\item AD:common:up\_axi:1.0.0
	\end{itemize}
\end{itemize}


\subsubsection{wishbone\_standard\_1553 Depenecies}
\begin{itemize}
\item dep
	\begin{itemize}
	\item AFRL:utility:helper:1.0.0
	\item AFRL:device:up\_1553:1.0.0
	\item AFRL:bus:up\_wishbone\_standard:1.0.0
	\end{itemize}
\end{itemize}


\subsubsection{up\_1553 Depenecies}
\begin{itemize}
\item dep
	\begin{itemize}
	\item AFRL:utility:helper:1.0.0
	\item AFRL:device\_converter:axis\_1553:1.0.0
	\item AFRL:buffer:fifo
	\end{itemize}
\end{itemize}


\subsection{In a Project}
\par
First, pick a core that matches the target bus in question. Then connect the BUS1553 core to that bus. Once this is complete the PMOD pins will need
to be routed so they match the PMOD1553 device. Please see the schematic of the PMOD1553 for electrical connection details. All I/O's are 3.3volt.

\section{Architecture}
\par
This core is made up of other cores that are documented in detail in there source. The cores this is made up of are the,
\begin{itemize}
  \item \textbf{axis\_1553\_encoder} Encodes data from the RX FIFO and sends it to the PMOD1553 (see core for documentation).
  \item \textbf{axis\_1553\_decoder} Decodes data from the PMOD1553 and sends it to the TX FIFO (see core for documentation).
  \item \textbf{fifo} Used for RX and TX FIFO instances. Set to 16 words buffer max (see core for documentation).
  \item \textbf{up\_axi} An AXI Lite to uP converter core (see core for documentation).
  \item \textbf{up\_wishbone\_standard} A wishbone standard to uP converter core (see core for documentation).
  \item \textbf{up\_1553} Takes uP bus and coverts it to interface with the RX/TX FIFOs and the encoder/decoder (see module documentation for information \ref{Module Documentation}).
\end{itemize}

For register documentation please see up\_1553 in \ref{Module Documentation}

\section{Building}

\par
The BUS1553 is written in Verilog 2001. It should synthesize in any modern FPGA software. The core comes as a fusesoc packaged core and can be included in any other core. Be sure to make sure you have meet the dependencies listed in the previous section. Linting is performed by verible using the lint target.

\subsection{fusesoc}
\par
Fusesoc is a system for building FPGA software without relying on the internal project management of the tool. Avoiding vendor lock in to Vivado or Quartus.
These cores, when included in a project, can be easily integrated and targets created based upon the end developer needs. The core by itself is not a part of
a system and should be integrated into a fusesoc based system. Simulations are setup to use fusesoc and are a part of its targets.

\subsection{Source Files}

\subsubsection{axi\_lite\_1553 File List}
\begin{itemize}
\item src
	\begin{itemize}
	\item src/axi\_lite\_1553.v
	\end{itemize}
\item tb\_cocotb
	\begin{itemize}
	\item {'tb/tb\_cocotb\_axi\_lite.py': {'file\_type': 'user', 'copyto': '.'}}
	\item {'tb/tb\_cocotb\_axi\_lite.v': {'file\_type': 'verilogSource'}}
	\end{itemize}
\end{itemize}


\subsubsection{wishbone\_standard\_1553 File List}
\begin{itemize}
\item src
	\begin{itemize}
	\item src/wishbone\_standard\_1553.v
	\end{itemize}
\item tb\_cocotb
	\begin{itemize}
	\item {'tb/tb\_cocotb\_wishbone\_standard.py': {'file\_type': 'user', 'copyto': '.'}}
	\item {'tb/tb\_cocotb\_wishbone\_standard.v': {'file\_type': 'verilogSource'}}
	\end{itemize}
\item tb
	\begin{itemize}
	\item tb/tb\_wishbone\_slave.v
	\end{itemize}
\end{itemize}


\subsubsection{up\_1553 File List}
\begin{itemize}
\item src
	\begin{itemize}
	\item[$\space$] Type: verilogSource
	\item src/up\_1553.v
	\end{itemize}
\item tb
	\begin{itemize}
	\item[$\space$] Type: verilogSource
	\item tb/tb\_up\_1553.v
	\end{itemize}
\end{itemize}


\subsection{Targets}

\subsubsection{axi\_lite\_1553 Targets}
\begin{itemize}
\item default
	\begin{itemize}
	\item[$\space$] Info: Default for IP intergration.
	\end{itemize}
\item lint
	\begin{itemize}
	\item[$\space$] Info: Lint with Verible
	\end{itemize}
\item sim\_cocotb
	\begin{itemize}
	\item[$\space$] Info: Cocotb unit tests
	\end{itemize}
\end{itemize}


\subsubsection{wishbone\_standard\_1553 Targets}
\begin{itemize}
\item default
	\begin{itemize}
	\item[$\space$] Info: Default for IP intergration.
	\end{itemize}
\item lint
	\begin{itemize}
	\item[$\space$] Info: Lint with Verible
	\end{itemize}
\item sim\_cocotb
	\begin{itemize}
	\item[$\space$] Info: Cocotb unit tests
	\end{itemize}
\end{itemize}


\subsubsection{up\_1553 Targets}
\begin{itemize}
\item default
	\begin{itemize}
	\item[$\space$] Info: Default for IP intergration.
	\item src
	\item dep
	\end{itemize}
\item sim
	\begin{itemize}
	\item[$\space$] Info: Base simulation using icarus as default.
	\item src
	\item dep
	\item tb
	\end{itemize}
\end{itemize}


\subsection{Directory Guide}

\par
Below highlights important folders from the root of the directory.

\begin{enumerate}
  \item \textbf{docs} Contains all documentation related to this project.
    \begin{itemize}
      \item \textbf{manual} Contains user manual and github page that are generated from the latex sources.
    \end{itemize}
  \item \textbf{src} Contains source files for the core
  \item \textbf{tb} Contains test bench files for iverilog and cocotb
    \begin{itemize}
      \item \textbf{cocotb} testbench files
    \end{itemize}
\end{enumerate}

\newpage

\section{Simulation}
\par
There are a few different simulations that can be run for this core.

\subsection{cocotb}
\par
Cocotb is the only method for simulating the various interations of the bus\_1553 core. At the moment there is a
axi\_lite, wishbone\_standard, and uP based versions. This is currently set to use icarus as the sim tool for cocotb.

\par
To run the wishbone sim use the command below.
\begin{lstlisting}[language=bash]
fusesoc run --target sim_cocotb AFRL:device:wishbone_standard_1553:1.0.0
\end{lstlisting}

\par
To run the axi\_lite sim use the command below.
\begin{lstlisting}[language=bash]
fusesoc run --target sim_cocotb AFRL:device:axi_lite_1553:1.0.0
\end{lstlisting}

\par
To run the uP sim use the command below.
\begin{lstlisting}[language=bash]
fusesoc run --target sim_cocotb AFRL:device:up_1553:1.0.0
\end{lstlisting}


\newpage

\section{Module Documentation} \label{Module Documentation}

\par
up\_1553 is the module that integrates the AXI streaming 1553 encoder/decoder.
This includes FIFO's that have there inputs/outputs for data tied to registers mapped in the uP bus.
The uP bus is the microprocessor bus based on Analog Devices design. It resembles a APB bus in design,
and is the bridge to other buses BUS1553 can use. This makes changing for AXI Lite, to Wishbone to whatever
quick and painless.

\par
axi\_lite\_1553 module adds a AXI Lite to uP (microprocessor) bus converter. The converter is
from Analog Devices.

\par
wishbone\_standard\_1553 module adds a Wishbone Standard to uP (microprocessor) bus converter. This
converter was designed for Wishbone Standard only, NOT pipelined.

\vspace{15mm}
\par
The next sections document these modules in great detail. up\_1553 contains the register map explained, and what the various bits do.

